\documentclass[11pt]{article}

\usepackage[margin=1in]{geometry}
\usepackage{amsmath,amssymb,amsthm,mathtools}
\usepackage{graphicx}
\usepackage[hidelinks]{hyperref}
\usepackage{enumitem}
\usepackage{microtype}

\newtheorem{theorem}{Theorem}
\newtheorem{lemma}[theorem]{Lemma}
\newtheorem{proposition}[theorem]{Proposition}
\newtheorem{corollary}[theorem]{Corollary}
\newtheorem{definition}[theorem]{Definition}
\theoremstyle{remark}
\newtheorem{remark}[theorem]{Remark}

\newcommand{\R}{\mathbb R}
\newcommand{\C}{\mathbb C}
\newcommand{\Hh}{\mathbb H}
\newcommand{\fg}{\mathfrak g}
\newcommand{\fh}{\mathfrak h}
\newcommand{\fs}{\mathfrak s}
\newcommand{\fz}{\mathfrak z}
\newcommand{\so}{\mathfrak{so}}
\newcommand{\su}{\mathfrak{su}}
\newcommand{\spalg}{\mathfrak{sp}}
\DeclareMathOperator{\SO}{SO}
\DeclareMathOperator{\SU}{SU}
\DeclareMathOperator{\U}{U}
\DeclareMathOperator{\Sp}{Sp}
\DeclareMathOperator{\Spin}{Spin}
\DeclareMathOperator{\Iso}{Iso}
\DeclareMathOperator{\End}{End}
\DeclareMathOperator{\Lie}{Lie}

\title{The Five-Dimensional Skew-Hermitian Lie Algebra Problem}
\author{Automatic Math Discovery Smoke Test}
\date{June 14, 2026}

\begin{document}
\maketitle

\section*{Source Problems}

The source paper is Kadets--Martin--Paya,
\emph{Recent progress and open questions on the numerical index of Banach
spaces} \cite{KMP}.  Problem 27 asks for the possible values of
\(\dim \mathcal Z(X)\), where \(\mathcal Z(X)\) is the Lie algebra of
skew-hermitian operators on a finite-dimensional real Banach space \(X\).
Problem 28 asks specifically for the five-dimensional case.  The source notes
that dimensions \(3\) and \(4\) were known and that the first open case is
dimension \(5\).

\begin{center}
\includegraphics[width=\textwidth]{figures/open_problem_crop.png}
\end{center}

\section*{Status of This Packet}

This packet is a candidate full solution of Problem 28, the five-dimensional
case.  It also records a promising but currently conditional route for
Problem 27 in all finite dimensions.

The previous promoted draft claimed the all-dimensional route as complete.
That was too strong.  Its deficit-summand theorem required a nonzero proper
\(G\)-invariant summand, which is false: for instance
\(\U(2)\leq \SO(4)\) acting on \(\C^2_\R\) has
\[
 \dim \U(2)=4,\qquad w(4)=5>4,
\]
but the real representation is irreducible.  Likewise
\(\SU(3)\curvearrowright \C^3_\R\) gives a strict triangular deficit without
a proper invariant summand.  The correct all-dimensional target must allow
the subspace \(W\) to be all of \(V\), and it must not require \(W\) to be
\(G\)-invariant.  What is needed for the Banach-space enlargement argument is
a subgroup fixing \(W^\perp\) pointwise and acting on \(W\) by a proper
sphere-transitive subgroup of \(\SO(W)\).

\section*{Proof Intuition for Problem 28}

Auerbach's theorem puts the connected linear isometry group of a
finite-dimensional real Banach space inside an orthogonal group, so
\(\mathcal Z(X)\) is a compact Lie subalgebra of \(\so(n)\).  Rosenthal's
gap theorem says that if \(\dim\mathcal Z(X)>(n-1)(n-2)/2\), then \(X\) is
Euclidean.  In dimension \(5\), this leaves only the possible obstruction
\(\dim\mathcal Z(X)=5\): values \(7,8,9\) are excluded by Rosenthal, and
\(10\) is the Euclidean case.

Thus the five-dimensional problem reduces to a small compact-Lie-algebra
check inside \(\so(5)\).  A compact \(5\)-dimensional Lie algebra is either
abelian of dimension \(5\), or has semisimple part \(\su(2)\) and a
two-dimensional center.  Neither embeds faithfully into \(\so(5)\): an
abelian subalgebra of \(\so(5)\) has rank at most \(2\), and the possible
\(\su(2)\)-module decompositions in dimension \(5\) leave centralizer rank at
most \(1\).  Hence dimension \(5\) cannot occur.  The remaining displayed
values are realized by finite \(\ell_\infty\)-sums of Euclidean blocks.

\section{The Five-Dimensional Classification}

\begin{theorem}[Problem 28]\label{thm:five}
If \(X\) is a five-dimensional real Banach space, then
\[
 \dim\mathcal Z(X)\in\{0,1,2,3,4,6,10\}.
\]
Conversely, every value in this set is attained by a five-dimensional real
Banach space.
\end{theorem}

\begin{proof}
By the finite-dimensional Auerbach reduction recalled in \cite{KMP}, after
choosing an auxiliary Euclidean structure on the underlying real vector space,
the identity component of the linear isometry group of \(X\) is a compact
subgroup of \(\SO(5)\), and its Lie algebra is \(\mathcal Z(X)\).  Hence
\(\mathcal Z(X)\) is a compact Lie subalgebra of \(\so(5)\).

Rosenthal's gap theorem for Banach-space isometry algebras
\cite{RosenthalLie}, as quoted in \cite{KMP}, implies that if
\[
 \dim\mathcal Z(X)>\frac{(5-1)(5-2)}2=6,
\]
then \(X\) is Euclidean and \(\dim\mathcal Z(X)=\dim\so(5)=10\).  Therefore
the only value not already eliminated among \(0,1,\ldots,6,10\) is \(5\).
It remains to show that no compact \(5\)-dimensional Lie subalgebra of
\(\so(5)\) can occur.

Let \(\fg\subseteq\so(5)\) be compact with \(\dim\fg=5\).  Its semisimple
part is either \(0\) or \(\su(2)\), since the smallest nonzero compact simple
Lie algebra has dimension \(3\).  If \(\fg\) is abelian, then it has
dimension at most the rank of \(\so(5)\), namely \(2\), contradiction.  Thus
\[
 \fg\cong \su(2)\oplus \R^2
\]
as a compact Lie algebra.

Consider the underlying faithful orthogonal representation on \(\R^5\) and
restrict it to the \(\su(2)\)-summand.  The nontrivial irreducible real
\(\su(2)\)-modules of dimension at most \(5\) have dimensions \(3,4,5\):
the \(3\)- and \(5\)-dimensional modules are of real type, and the
\(4\)-dimensional module is of quaternionic type.

If the restriction is irreducible of dimension \(5\), the skew centralizer
in \(\so(5)\) is \(0\), so no two-dimensional center can act faithfully.  If
the decomposition is \(4+1\), the \(4\)-dimensional quaternionic block has
orthogonal centralizer \(\spalg(1)\), whose maximal abelian subalgebras are
one-dimensional, and the trivial line contributes no skew operator.  Again a
two-dimensional center cannot inject.  If a \(3\)-dimensional irreducible
block occurs, the remaining fixed part has dimension at most \(2\), whose
skew algebra has rank at most \(1\); the \(3\)-dimensional real-type block
has zero skew centralizer.  This also prevents a faithful two-dimensional
central action.  These cases exhaust dimension \(5\), so
\(\dim\mathcal Z(X)\neq5\).

The listed values are realized as follows, where \(H_k\) is \(k\)-dimensional
Euclidean space and the sums are finite \(\ell_\infty\)-sums:
\[
\begin{array}{c|c}
\dim\mathcal Z(X) & X\\ \hline
0 & \ell_\infty^5\\
1 & H_2\oplus_\infty \ell_\infty^3\\
2 & H_2\oplus_\infty H_2\oplus_\infty\R\\
3 & H_3\oplus_\infty \ell_\infty^2\\
4 & H_3\oplus_\infty H_2\\
6 & H_4\oplus_\infty \R\\
10 & H_5.
\end{array}
\]
For such a sum the identity component of the isometry group is the product
of the orthogonal groups on the Euclidean blocks, so the Lie algebra dimension
is the displayed sum of \(\binom{k}{2}\)'s.
\end{proof}

\section{Triangular Bookkeeping for Problem 27}

For \(r\geq 2\), put \(T_r=\binom r2\), and define
\[
 w(N)=\min\left\{\sum_i r_i:
       N=\sum_i \binom{r_i}{2},\ r_i\geq 2\right\},
 \qquad w(0)=0.
\]
Because \(T_2=1\), \(w(N)\) is defined for every \(N\geq0\).  Let
\[
 P(n)=
 \left\{\sum_i \binom{k_i}{2}:\sum_i k_i=n,\ k_i\geq1\right\}.
\]
Then \(N\in P(n)\) if and only if \(w(N)\leq n\).

\begin{lemma}[elementary support estimates]\label{lem:w-estimates}
The following estimates hold.
\begin{enumerate}[label=\textup{(\roman*)}]
\item \(w(A+B)\leq w(A)+w(B)\).
\item \(w(\binom r2)=r\).
\item For \(k\geq3\), \(w(k^2-1)\leq4k-4\).
\item For \(k\geq2\), \(w(k^2)\leq2k+1\).
\item For \(k\geq1\), \(w(k(2k+1))\leq2k+1\).
\end{enumerate}
\end{lemma}

\begin{proof}
Subadditivity follows by concatenating triangular decompositions.  The
identity \(w(\binom r2)=r\) follows from
\(\binom a2+\binom b2<\binom{a+b}2\) for \(a,b\geq2\): with total support
\(S<r\), the largest triangular sum is at most \(\binom S2<\binom r2\).
For (iii), use
\[
 k^2-1=\binom{k+1}{2}+\binom{k-1}{2}+(k-2)\binom22.
\]
For (iv), use \(k^2=\binom{k+1}{2}+\binom k2\).  Finally
\[
 k(2k+1)=\binom{2k+1}{2}.
\]
\end{proof}

\begin{proposition}[realization in every dimension]\label{prop:realization}
Every value in \(P(n)\) is attained as \(\dim\mathcal Z(X)\) for some
\(n\)-dimensional real Banach space \(X\).
\end{proposition}

\begin{proof}
If \(n=\sum_j k_j\), take
\[
 X=H_{k_1}\oplus_\infty\cdots\oplus_\infty H_{k_r}.
\]
The identity component of the isometry group preserves the finite
\(\ell_\infty\)-summand decomposition and acts as the product of orthogonal
groups on the Euclidean factors.  Hence
\[
 \dim\mathcal Z(X)=\sum_j\dim\so(k_j)=\sum_j\binom{k_j}{2}.
\]
\end{proof}

\section{Conditional Route for Problem 27}

The expected all-dimensional answer to Problem 27 is:
\[
 \dim\mathcal Z(X)\in P(n)
 \quad\text{for every \(n\)-dimensional real Banach space \(X\),}
\]
with the converse supplied by Proposition~\ref{prop:realization}.  The
current packet still does not claim this as proved, but the structural flaw in
the earlier proof has been repaired: the right object is a defective subspace,
not a proper \(G\)-invariant summand.  The remaining proof-grade work is the
compact representation input described below.

\subsection*{Conditional Dependencies}

\paragraph{Dependency A: a proof-grade small-module classification.}
One needs the following representation-theoretic statement, split into
machine-checkable pieces rather than a table summary.

\begin{itemize}[leftmargin=2em]
\item If \(\fs\) is compact simple and \(E\) is an irreducible faithful real
\(\fs\)-module with \(w(\dim\fs)>\dim_\R E\), then
\[
(\fs,E)=(\su(k),\C^k_\R)\quad(k\geq3)
\quad\text{or}\quad
(\mathfrak g_2,\R^7),
\]
up to finite covers, duals, and low-rank isomorphisms.
\item If \(\fs\oplus\R\) acts irreducibly and effectively with
\(w(\dim\fs+1)>\dim_\R E\), then
\[
(\fs\oplus\R,E)=(\mathfrak u(k),\C^k_\R),
\]
again up to the standard low-rank identifications.
\item If
\(\fh=\fz\oplus\fs_1\oplus\cdots\oplus\fs_r\) with \(r\geq2\) acts
irreducibly and effectively, then
\[
 w(\dim\fh)\leq\dim_\R E.
\]
\end{itemize}

The supplied supporting file
\path{supporting_small_module_lemma.tex} gives a substantial proof candidate.
The Weyl-dimension and product scans described in the verification notes found
no suspicious cases.  Still, this is the remaining bottleneck: the
low-degree tables, real/quaternionic type bookkeeping, and product-factor
case should be expanded or cited precisely before Problem 27 is marked full.

\paragraph{Dependency B: compact Lie algebra splitting over semisimple ideals.}
This part is standard and proof-grade:
if
\[
0\longrightarrow\mathfrak k\longrightarrow\mathfrak q
 \longrightarrow\mathfrak a\longrightarrow0
\]
is an exact sequence of compact Lie algebras and
\(\fs\subseteq\mathfrak a\) is a semisimple ideal, then there is a
semisimple ideal \(\widetilde\fs\subseteq\mathfrak q\) mapping
isomorphically onto \(\fs\), and
\[
[\widetilde\fs,\mathfrak k]=0.
\]
Indeed, a compact Lie algebra is the direct sum of its center and compact
simple ideals.  The simple ideals of \(\mathfrak q\) that map nontrivially
onto the simple factors of \(\fs\) map isomorphically, and the remaining
kernel ideals commute with them.

\paragraph{Dependency C: centralizer coupling.}
The reducible step uses the following centralizer input.  Let \(A\) be one of
the sphere-transitive atoms from Dependency A, acting on a real block \(E\).
Let \(S\) be its transitive semisimple factor.  Suppose \(S\) also acts
nontrivially on another irreducible real summand
\[
 F=U\otimes_DD^m,\qquad D=\End_S(U)\in\{\R,\C,\Hh\}.
\]
If \(Q\) acts on \(E\oplus F\), projects onto \(A\) on \(E\), and has kernel
over \(A\) contained in the \(S\)-centralizer on \(F\), then
\[
 w(\dim Q)\leq \dim E+\dim_\R F.
\]
The appendix and the centralizer-coupling interval attempt note give the
interval proof.  The arithmetic was also stress-tested by the centralizer
verifier scripts.  The only representation-theoretic reduction used there is
the same minimal-module input included in Dependency A.

\subsection*{Corrected Deficit-Subspace Theorem}

For \(G\leq\SO(V)\) compact and connected, with Lie algebra \(\fg\), and for
a nonzero orthogonal subspace \(W\leq V\), define
\[
 \fg[W]=\{A\in\fg:\ A|_{W^\perp}=0\}.
\]
Since \(A\) is skew-symmetric, \(A|_{W^\perp}=0\) implies \(A(W)\subseteq W\).

\begin{definition}
The subspace \(W\) is \emph{defective} for \(G\) if \(\fg[W]\) contains a
compact connected subalgebra acting transitively on the Euclidean unit sphere
of \(W\), but
\[
 \so(W)\not\subseteq \fg[W].
\]
\end{definition}

\begin{theorem}[conditional corrected deficit theorem]\label{thm:defective}
Assume Dependencies A--C.  Let \(G\leq\SO(V)\) be compact, connected, and
effective.  If
\[
 w(\dim G)>\dim V,
\]
then \(V\) contains a nonzero defective subspace \(W\).
\end{theorem}

\begin{proof}
Assume not, and choose a counterexample with \(\dim V\) minimal.  Decompose
\[
V=E_1\oplus\cdots\oplus E_r
\]
into irreducible real \(G\)-summands.  Let \(H_i\) be the image of \(G\) on
\(E_i\), and let \(K_i\) be the kernel of the projection \(G\to H_i\).  At
the Lie algebra level,
\[
\dim G=\dim K_i+\dim H_i.
\]

Suppose \(w(\dim H_i)\leq\dim E_i\) for some \(i\).  Then \(K_i\) acts
effectively on \(E_i^\perp\).  If \(K_i\) had a defective subspace
\(W\leq E_i^\perp\), then the same \(W\) would be defective for \(G\):
the transitive subalgebra still fixes the full orthogonal complement of
\(W\) in \(V\), and if \(\so(W)\subseteq\fg[W]\), those block rotations
would act trivially on \(E_i\) and hence would already lie in the
corresponding kernel algebra.  Therefore \(K_i\) has no defective subspace on
\(E_i^\perp\).  By minimality,
\[
w(\dim K_i)\leq\dim E_i^\perp.
\]
Subadditivity gives
\[
w(\dim G)\leq w(\dim H_i)+w(\dim K_i)
\leq \dim E_i+\dim E_i^\perp=\dim V,
\]
a contradiction.  Hence every irreducible image satisfies
\[
w(\dim H_i)>\dim E_i.
\]

By Dependency A, every \(H_i\) is one of the strict-deficit
sphere-transitive atoms \(SU(k)\), \(U(k)\), or \(G_2\), up to the listed
standard identifications.  If \(r=1\), then \(W=V\) is defective: the
semisimple transitive factor acts transitively on the sphere, while the full
\(\so(V)\) is not contained in the proper atom algebra.  Thus \(r\geq2\).

Fix \(E=E_i\), put \(A=H_i\), and let \(S\) be the transitive semisimple
factor of \(A\).  By Dependency B, the \(S\)-factor lifts to a semisimple
ideal \(\widetilde S\subseteq\fg\) commuting with the kernel over \(A\).  If
\(\widetilde S\) acts trivially on \(E^\perp\), then \(E\) is defective,
again because the image of \(\fg[E]\) on \(E\) lies inside the proper atom
algebra \(A<\SO(E)\).  Hence \(\widetilde S\) acts nontrivially on some
other irreducible summand \(F\).

Let \(Q\) be the image of \(G\) on \(E\oplus F\), and let \(N\) be the kernel
of the projection \(G\to Q\).  The projection \(Q\to A\) has kernel contained
in the \(\widetilde S\)-centralizer on \(F\), so Dependency C gives
\[
w(\dim Q)\leq \dim E+\dim F.
\]
The kernel \(N\) acts effectively on \((E\oplus F)^\perp\).  As above, a
defective subspace for \(N\) there would also be defective for \(G\), so
minimality gives
\[
w(\dim N)\leq\dim(E\oplus F)^\perp.
\]
Thus
\[
w(\dim G)\leq w(\dim Q)+w(\dim N)
\leq \dim E+\dim F+\dim(E\oplus F)^\perp=\dim V,
\]
the final contradiction.
\end{proof}

\subsection*{How the Dependencies Would Prove Problem 27}

Assume Dependencies A--C.  Let \(G=\Iso_0(X)\).  By the same Auerbach
reduction used above,
\[
G\leq\SO(V),\qquad \Lie(G)=\mathcal Z(X),
\]
for an auxiliary Euclidean structure on the underlying real vector space
\(V\) of \(X\).  If
\[
w(\dim\mathcal Z(X))>\dim X,
\]
Theorem~\ref{thm:defective} gives a nonzero defective subspace \(W\).  For
each \(u\in W^\perp\), the function
\[
x\mapsto \|u+x\|_X\qquad(x\in W)
\]
is invariant under a connected subgroup acting transitively on the Euclidean
unit sphere of \(W\).  Hence this slice norm is radial on \(W\), so the whole
norm is invariant under the full \(\SO(W)\) acting on \(W\) and fixing
\(W^\perp\).  But defectiveness says \(\so(W)\not\subseteq\fg[W]\), so this
strictly enlarges the connected isometry algebra, contradicting
\(G=\Iso_0(X)\).  Thus \(w(\dim\mathcal Z(X))\leq n\), which is equivalent
to \(\dim\mathcal Z(X)\in P(n)\).

\section*{Verification Notes}

The script \texttt{code/check\_dimension\_values.py} records the triangular
sets
\[
P(5)=\{0,1,2,3,4,6,10\},
\]
\[
P(6)=\{0,1,2,3,4,6,7,10,15\},
\]
and
\[
P(7)=\{0,1,2,3,4,5,6,7,9,10,11,15,21\}.
\]
These computations verify the finite bookkeeping and construction values,
not the conditional all-dimensional representation-theoretic dependencies.

The internal appendix scripts retain the finite and symbolic checks for the
centralizer-coupling bounds and the low-dimensional \(n=6,7\) bookkeeping.
The centralizer interval scripts report no failures in the tested ranges, and
the Weyl/product verification aids report no suspicious small-module or
product-tensor rows.  These computations support Dependencies A and C, but
they do not replace a written citation or line-by-line proof of the
low-degree compact-representation tables in Dependency A.

\section*{Novelty Check}

Earlier bounded local searches over the parsed corpus and run artifacts found
no answer to Problem 28 beyond the source statement.  Bounded web searches on
June 13, 2026 for exact phrases around \(\dim Z(X)\), finite-dimensional real
Banach spaces, skew-hermitian operators, and Rosenthal's Lie-algebra paper
found adjacent numerical-index and skew-hermitian-operator literature, but no
direct later answer to Problems 27 or 28.  The Problem 28 result should still
be treated as a candidate pending expert review.

\section*{Human Review Focus}

For Problem 28, review the compact Lie algebra exclusion of
\(\dim\mathcal Z(X)=5\), especially the centralizer-rank check for the
\(\su(2)\)-module decompositions.

For the all-dimensional Problem 27 route, do not yet treat this packet as a
full proof.  The corrected defective-subspace architecture is now included;
the main review targets are:
\begin{enumerate}
\item expand or precisely cite the low-degree compact-representation table in
Dependency A, especially product tensors and real/quaternionic type;
\item check the minimal-module reduction used in the centralizer coupling
lemma;
\item verify the corrected defective-subspace induction in
Theorem~\ref{thm:defective};
\item ensure the final Banach-space enlargement is strict through the
definition of defectiveness.
\end{enumerate}

\begin{thebibliography}{9}

\bibitem{BorelSpheres}
A. Borel,
\emph{Le plan projectif des octaves et les sph\`eres comme espaces homog\`enes},
C. R. Acad. Sci. Paris \textbf{230} (1950), 1378--1380.

\bibitem{BourbakiLie}
N. Bourbaki,
\emph{Lie Groups and Lie Algebras, Chapters 4--6},
Springer, 2002.

\bibitem{KMP}
V. Kadets, M. Mart\'in, and R. Pay\'a,
\emph{Recent progress and open questions on the numerical index of Banach
spaces},
arXiv:math/0605781.

\bibitem{RosenthalLie}
H. P. Rosenthal,
\emph{The Lie algebra of a Banach space},
in: Banach spaces (Columbia, Mo., 1984), Lecture Notes in Math.
\textbf{1166}, Springer, 1985, 129--157.

\end{thebibliography}

\end{document}
